% !TEX root = ../jlt.tex
\section{Introduction}
\label{sec:intro}

The case for functional programming has been made the same way for decades, and it is a good case; even a compiler's own intermediate form, static single assignment, is a functional language~\cite{appel1998ssa}.
A function that depends on its arguments alone can be replaced by its value: programs can be read as equations, rewritten by them, and reasoned about by them~\cite{backus1978can}.
A function that is polymorphic in a type cannot inspect that type, so its type alone says a great deal about what it does~\cite{wadler1989theorems}.
A class of types with operations and laws lets generic code be written once against the operations and reasoned about by the laws~\cite{wadler1989adhoc}.
The two are one mechanism with two sources of equations: a free theorem is proved from the type alone, for every term, and a law is declared of an instance.
Either becomes a rewrite by Leibniz's law, which Rule~1 conserves (Section~\ref{sec:posture}).
For the declared laws Birkhoff's theorem makes the reasoning complete: an equation follows from the laws by equational reasoning exactly when it holds in every model of them~\cite{birkhoff1935structure,burris1981universalalgebra}.
A class with laws is an equational theory and its instances the models, a variety, and a stronger class has fewer instances, so the constraints describe the structure and the structure licenses the optimisation.
The goal is correctness by construction in Dijkstra's sense~\cite{dijkstra1976discipline}.
It is usually made in a language built around it, sometimes experimental or niche.
The underlying concepts, however, are often older, run deeper and have broader reach than is assumed:
the most widely written functional language may well be SQL's \texttt{SELECT}, a declarative query over relations~\cite{iso9075sql}, whose algebra is Codd's~\cite{codd1970relational}, with Tarski's calculus of relations~\cite{tarski1941calculus} its older relative.

This paper makes the case in a mainstream multi-paradigm systems language: \Cpp{23}.
Much of machine learning and most embedded devices run on \Cpp{}, and a growing number of the people who write it arrive from Python with a functional vocabulary.
Ubiquity and performance are why one comes to such a language in the first place: a compiler for nearly every machine, and abstraction without overhead by its zero-overhead principle~\cite{stroustrup2005abstraction}.\footnote{Stroustrup has since called the phrase an understatement: ``It's compiled away. This is why I talk about zero overhead abstraction, and people are beginning to take me to task for that because that's understating and underestimating the ability of the \Cpp{} compiler. We can do negative overhead abstraction''~\cite{stroustrup2026peterman}.
Listing~\ref{lst:sites} is a case: a count over all sequences compiles away to a constant, a matcher's construction to its table, and a law with a length bound to a vectorised scan.}
% Should probably go to conclusions or something:
%For them the repository binds the library's combinators into Python as handles to the compiled core, so the laws declared in \Cpp{} hold in an interactive session: the specialisations they license run there at run time, and only what the compiler evaluates of them is fixed when the bindings are built (Section~\ref{sec:binding}).
A functional fragment of \Cpp{23} is worth drawing because the whole may be too large to work in with confidence.
Such a fragment is worth drawing now because \Cpp{23} is the first standard with every ingredient: compile-time evaluation of ordinary functions and of most of the standard library, concepts, and folds over ranges.
Compliance with such a fragment is also easier to check than a decade ago: for example, a language model of code~\cite{chen2021evaluating} can be instructed to proof-read every change against its rules.
Such a review is a read, not a proof (Section~\ref{sec:limits}).

Machine-checked mathematics and mainstream software sit at opposite corners.
A proof checked in Lean~\cite{demoura2021lean} is a white box, true on the checker's word however it was found~\cite{hubert2025olympiad}, and whether it is interesting is left to a human reader~\cite{rudin2019stop,tao2026palomar}.
A program in production is often the reverse: interesting, even commercially viable, from the start, and as opaque as the search that may have written it~\cite{turpin2023unfaithful,lindsey2025biology}; the open question is whether it is true, whether it behaves as its authors intended.
Yet a program is a formal artefact too: a compiler translates it towards a formalised machine, so every claim about it is a claim about something formal, and a compile-time check is such a claim, which the compiler verifies.
For generic code, \Cpp{} lacks just such a check.
\cppinline{concept}s as proposed for \Cpp{0x} checked a template's definition against its constraints~\cite{gregor2006concepts,dosreis2006specifying}, but as standardised in \Cpp{20} they check only its uses~\cite{stroustrup2020thriving}, so a template's body may use operations its constraints do not name.
The designers named the problem at the outset, in 2003: ``crucial aspects of the programmer's understanding of a template cannot be expressed directly in the code for compilers and other tools to use''~\cite{stroustrup2003concepts}.

\Cpp{23} therefore contains the material for a typed functional fragment but does not enforce its boundary, and this paper documents a discipline that does, recovering some of what a functional language guarantees at the cost of three rules.
The cost is paid once, in how a program is written, and what it buys is paid back at every use: a law holds in every model of it, a specialisation it licenses pays on every call, and a checked statement serves every reader.
The fragment characterised here is standard \Cpp{23}, restricted, not extended: it complies with the principles of the host language, and is accepted by any conforming compiler.
The compiler's optimisation pipelines can be brought to bear:
an optimising compiler already reasons about a program, by abstract interpretation~\cite{cousot1977abstract} and by algebraic laws it is told it may assume, as \texttt{-ffast-math} lets it for floating point, and which in Jlt are declared.

The fragment is the \emph{Juliet Posture}, or Jlt, and it has two boundaries, both drawn here.
The outer boundary is drawn by the three rules: Jlt is the proper subset of \Cpp{23} they define.
We call it a posture because it asks only for a style of programming: every ingredient is in the ISO standard~\cite{cpp23standard}, implemented by toolchains in wide use, and what Jlt adds is the programmer's discipline.
The compiler holds some of the rules and the programmer the rest, and a linter helps but is not required (Table~\ref{tab:posture}).
The inner boundary is drawn by the translation: the \emph{kernel}, Featherweight Jlt (FJlt), is the part of Jlt on which a translation to System~F ($\llbracket\cdot\rrbracket$) is defined, and it inherits from System~F termination, parametricity and a value.
Between the two lies \emph{sugar}, in Landin's word~\cite{landin1964mechanical}, which this paper names and characterises without proof (Appendix~\ref{app:fjlt:closure}), the arrangement being the Haskell Report's, a kernel and the translations into it~\cite{marlow2010haskell}.

Jlt is thus an embedded domain-specific language in Hudak's sense~\cite{hudak1996building}, whose domain is System~F and whose host is \Cpp{23}.
Its programs are ordinary \Cpp{23} programs, type-checked and compiled by an unmodified toolchain, Clang~23 in this case (Section~\ref{sec:limits}).
It is not a language this paper designs and develops as an architect might develop a house.
Instead it characterises a fragment of a pre-existing ecosystem, as a traveller might discover a city.
The claim here is one of existence: \Cpp{23} contains an embedding of a sizeable fragment of System~F, and the listings exhibit it.
No claim is made that Jlt is unique or canonical, although it seems likely that another fragment drawn to the same ends would look much like it.
As in the tagless-final style of Carette, Kiselyov and Shan~\cite{carette2009finally}, the embedded terms are generic code over an interface, here a \cppinline{concept}, and their interpretations are its models.
The posture's name is a nod to structural typing~\cite{pierce2002tapl,cardelli1996typesystems}: in \Cpp{23} a \cppinline{concept} is a predicate on a type's structure, not its lineage, and so a rose by any other name would smell as sweet~\cite{shakespeare1597romeo}.
Figure~\ref{fig:arrows} sketches the arrows between these pieces.

\begin{figure}[tb]
\centering
\begin{tikzcd}[column sep=2.4em, row sep=2.2em]
  & \text{Jlt} \arrow[dr, hook, "{-}\,\text{rules}"] \arrow[dl, dashed, bend right=25, "{-}\,\text{sugar}"'] \\[-0.8em]
\text{FJlt} \arrow[ur, hook, bend right=25, "{+}\,\text{sugar}"'] \arrow[d, bend right=25, "\llbracket\cdot\rrbracket"']
  &
  & \text{\Cpp{23}} \arrow[d, "\texttt{clang++}"] \\
\text{System~F} \arrow[u, hook, bend right=25, "\text{embed}"'] \arrow[d, "\text{eval}"']
  &
  & \text{LLVM IR} \arrow[d, "\text{run}"] \\
\text{value} \arrow[rr, equal]
  &
  & \text{value}
\end{tikzcd}
\caption{The arrows of the paper.
The inclusions of FJlt in Jlt and of Jlt in \Cpp{23} are proper, and each is labelled with what it adds: Jlt is the kernel and the sugar, \Cpp{23} is Jlt without the rules.
The translation $\llbracket\cdot\rrbracket$ is total on FJlt, and $\mathit{embed}$, defined on its image, a non-empty part of System~F, is a section of it~\cite{nlab:retract}.
Its composite with the two inclusions is a partial map from System~F, the embedding that makes Jlt an embedded language.
\texttt{clang++} is defined on every \Cpp{23} program within the implementation's limits.
Solid arrows are established: for every closed $t$ of base type in the image, $\mathit{eval}\,t = \mathit{run}\,(\texttt{clang++}\,(\mathit{embed}\,t))$.
The dashed arrow is a working hypothesis: $\mathit{desugar}$ is total on Jlt and a retraction of the inclusion.
$\mathit{desugar} \circ \mathit{incl} = \mathit{id}$ holds exactly, by construction.
$\mathit{incl} \circ \mathit{desugar} = \mathit{id}$ holds up to the value a program computes.
The bottom edge then holds for every Jlt program $p$: $\mathit{eval}\,\llbracket \mathit{desugar}\,p \rrbracket = \mathit{run}\,(\texttt{clang++}\,p)$.}
\label{fig:arrows}
\end{figure}

\paragraph{Contributions.}
The contribution is a map from a known discipline onto existing mainstream functionality, not a piece of software.
\begin{enumerate}
\item The posture: three rules on standard \Cpp{23}, each with what the compiler checks of it and what it leaves to the author (Section~\ref{sec:posture}).
\item Featherweight Jlt, the kernel: a core calculus whose syntax is a subset of \Cpp{23} as written, with typing, evaluation and a translation into System~F. Three theorems about the translation and three corollaries, with proof sketches (Section~\ref{sec:guarantees}, Appendix~\ref{app:fjlt:theorems}).
\item An adequacy argument from \Cpp{} to the calculus, which states its gaps rather than writing around them (Sections~\ref{sec:translation:adequacy} and~\ref{sec:limits}), and the sugar, named with its desugaring (Appendix~\ref{app:fjlt:closure}).
\item An observation about \Cpp{20} concepts: with overlapping constrained overloads, a generic function's behaviour can depend on a constraint it never declared, which is GHC's incoherent instance resolution, by default and silently.
Jlt adopts GHC's rule, and the repository's linter reports a call at a value of parameter type that breaks it (Appendix~\ref{app:fjlt:typing}, Section~\ref{sec:limits}).
\item The exhibits: every listing in this paper compiles and runs under a stock compiler at both binding times, by a test in the accompanying repository~\cite{Krautler_dedekind_2026}, and is held to the three rules by a mandatory review against them and by a mechanical linter in the same repository: that is what compliance means here.
\end{enumerate}
Not contributed: the translation, which is Wadler and Blott's~\cite{wadler1989adhoc}; the calculus-with-concepts, which is Siek and Lumsdaine's $F^G$~\cite{siek2005essential}; and complete definition checking: Section~\ref{sec:limits} says what the repository's linter reads and what it misses.

\paragraph{How to read this paper.}
The paper is written for authors of generic \Cpp{} libraries first, with System~F as the warranty and a rule for their overloads; for those who audit \Cpp{} against a policy and need to know what a tool checks; for readers of programming languages in the line of concepts and type classes; and for researchers who arrive from Python and want the performance of \Cpp{} with a guard rail.
The sections follow the arrows of Figure~\ref{fig:arrows}, one each.
Section~\ref{sec:posture} is the inclusion of Jlt in \Cpp{23}, the rules that carve Jlt out of it; Section~\ref{sec:translation:adequacy} is the inclusion of FJlt in Jlt and the dashed arrow back, which Jlt programs are kernel programs; Section~\ref{sec:translation} is $\llbracket\cdot\rrbracket$, the translation of Featherweight Jlt into System~F; Section~\ref{sec:guarantees} is \emph{eval}, the theorems read back from System~F; Section~\ref{sec:binding} is \texttt{clang++} and \emph{run}, one program at both binding times, and how instantiation relates to the translation.
The listings follow one example through the paper, and Figure~\ref{fig:structures} sketches the recursion it runs on: \cppinline{concat}, a monoid at the lengths that need no recursion, then the fold over the chain of a count under an arrow, then Fibonacci's matrix chain, which two declared laws fold in logarithmically many steps, then restriction sites, where the compiler partially evaluates a matcher in its pattern.
Section~\ref{sec:limits} says where the arrows are partial, Section~\ref{sec:related} where the work sits, and the calculus is in the appendices.

